\subsection{How Should I Act When My Energy Is Low?}

If your volitional energy is low, do not ask for a heroic restart. Reduce the
size of the target state until it can still win against framework inertia, and
use the smallest decision window you can reliably enter. The business goal is
not to prove strength; it is to preserve movement. In practice, that means
choosing one concrete action that is cheap enough to start, visible enough to
verify, and stable enough to create a foothold for the next step.

\begin{enumerate}[nosep]
  \item Name the target state in one sentence, but shrink it to the smallest
        useful version. If ``finish the report'' is too heavy, use ``open the
        file and write the first three lines.''
  \item Remove one obvious barrier before you try to start. Move the phone,
        open the document, place the book on the desk, or prepare the tools so
        the current state has less support.
  \item Choose the first action that is visible and almost impossible to
        misinterpret. A weak start is still better than a vague intention.
  \item Enter the decision window immediately. Low energy makes delay more
        expensive, because the old state regains control quickly.
  \item Stop after the first stable foothold if necessary. The purpose is to
        create a new trajectory, not to extract maximum output from a tired
        system.
  \item Repeat the same entry pattern next time so the transition becomes
        cheaper. Small history effects matter more than one dramatic effort.
\end{enumerate}

\begin{remark}[Low-energy rule]
When energy is low, the right question is not ``How do I finish everything?''
It is ``What is the smallest target state that can beat the current state
today?'' That question keeps the reader inside the regime of actionable
control, where barrier reduction and timing matter more than intensity.
\end{remark}

\begin{example}[A tired evening]
A reader comes home exhausted and wants to study. Instead of planning a full
two-hour session, they define the target state as ``read one page and write
one question.'' They place the book on the desk before dinner, silence the
phone, and begin during the first quiet minute after eating. Even if the
session ends after ten minutes, the transition succeeded because the new frame
was started while the old one was still weak.
\end{example}
